\subsection{The actual binomial deformation and its strict coefficient maps}

\paragraph{A-S09-15 (Proved; exact specialization maps and a quadratic lift obstruction).}
Let $B=\operatorname{Bin}(\mathbb Z)$, with its filtration
$F_nB$ by polynomial degree.  Write
\[
 B=\bigoplus_{n\geq0}\mathbb Z\binom{x}{n}
       \subset\mathbb Q[x],\qquad
 \mathcal R=\bigoplus_{n\geq0}h^nF_nB\subset B[h],
 \qquad b_n=h^n\binom{x}{n},\quad v=b_1.
 \tag{S09.44}
\]
Here $h,v$ have degree one and $b_n$ has degree $n$.
The following are actual algebra maps:
\[
 \mathcal R\longrightarrow\mathcal R/(h)
        \xrightarrow{\ \cong\ }\Gamma_{\mathbb Z}(v),
 \qquad
 \mathcal R\longrightarrow\mathcal R/(h-1)
        \xrightarrow{\ \cong\ }B.
 \tag{S09.45}
\]
The first sends $b_n$ to $\gamma_n(v)$; the second
sends it to $\binom{x}{n}$.  The algebra $\mathcal R$
is free as a $\mathbb Z[h]$-module.  Nevertheless:

\begin{enumerate}
\item No ideal of $\mathcal R$ containing $v$ admits
divided powers.  In particular the first specialization
in (S09.45) does not come from a strict PD algebra
structure on its source with $v$ in the PD ideal.
\item Put $S=\Gamma_{\mathbb Z}(T)$, with $T$ of degree
one.  There is no graded algebra map
$\mathcal R\to S[h]$ fixing $h$ and inducing the
map $v\mapsto T$ after reduction modulo $h$.
\item There is no algebra map $B\to S$ sending $x$ to
$T$.  Thus even this specified one-generator target
map cannot be transported from the divided-power
fiber to the binomial fiber by those specializations.
\end{enumerate}

\begin{proof}
The binomial basis in (S09.44) gives
\[
                     \mathcal R=
                 \bigoplus_{n\geq0}\mathbb Z[h]b_n.
 \tag{S09.46}
\]
Indeed $h^m\binom{x}{n}=h^{m-n}b_n$ for $m\geq n$,
and independence follows inside $\mathbb Q[x,h]$.
The exact multiplication formula is
\[
 b_m b_n=
 \sum_{j=0}^{\min(m,n)}
 \frac{(m+n-j)!}{j!(m-j)!(n-j)!}\,
                  h^j b_{m+n-j}.
 \tag{S09.47}
\]
It follows, for example, by counting pairs of subsets
of sizes $m,n$ according to the size $j$ of their
intersection, and hence is an identity of integer-valued
polynomials.  Modulo $h$, only $j=0$ remains, giving
the divided-power multiplication.  Modulo $h-1$ the
original binomial multiplication is recovered.
Formula (S09.46) proves that these are isomorphisms
onto the asserted fibers and proves flatness over
$\mathbb Z[h]$.

The case $m=n=1$ gives
\[
                         v^2=2b_2+h v.              \tag{S09.48}
\]
By (S09.46), the element on the right is not divisible
by $2$ in $\mathcal R$: its coefficient of the basis
element $b_1$ is $h$.  The PD identity
$v^2=2\gamma_2(v)$ is consequently impossible.  This
proves the first assertion without assuming any
particular formula for the proposed divided powers.

For the second assertion, a graded map as stated
necessarily sends $v$ to $T+kh$ for an integer $k$,
since the degree-one part of $S[h]$ is
$\mathbb ZT\oplus\mathbb Zh$.  Applying the map to
(S09.48) would give
\[
 2\,\operatorname{image}(b_2)=
  2T^{[2]}+(2k-1)hT+(k^2-k)h^2.                     \tag{S09.49}
\]
The three displayed degree-two basis elements of
$S[h]$ are independent, and $2k-1$ is odd.  Thus
the right side is not divisible by $2$.
Finally, specializing (S09.48) at $h=1$ shows that
a map in the third assertion would require
$2\operatorname{image}(\binom{x}{2})=2T^{[2]}-T$,
which is again impossible in the free abelian group
$S$.  This proves all three assertions.
\end{proof}

These calculations also apply to the simplex-zero
coefficient algebra of the standard admissible
PD de Rham theory $G$, since $G_0=S$.
They exclude the specified strict coefficient lifts
and PD structures; they do not exclude replacements
of the Rees family or other comparisons.

\paragraph{The proved derived comparison used in this audit.}
Kubrak--Shuklin--Zakharov define the filtered derived
binomial monad by nonabelian derived extension of
$F_nB$; their Remark~2.52 gives its actual associated
graded functors
\[
 \operatorname{gr}_n\mathbf L\operatorname{Bin}
                  \cong\mathbf L\Gamma^n.
 \tag{S09.50}
\]
For a perfect complex $M$ in cohomological degrees
$\leq-1$, the counit
$C_*^{\rm sing}(K(M),\mathbb Z)\to M$ dualizes to
$M^\vee\to C^*_{\rm sing}(K(M),\mathbb Z)$.
The free derived binomial adjunction extends this
specific map to
\[
 \mathbf L\operatorname{Bin}(M^\vee)
        \longrightarrow C^*_{\rm sing}(K(M),\mathbb Z).
 \tag{S09.51}
\]
Their Proposition~3.6 proves that (S09.51) is an
equivalence.  Its base case $M=\mathbb Z[1]$ uses
the d\'ecalage calculation
$\mathbf L\Gamma^n(\mathbb Z[-1])
\cong(\Lambda^n\mathbb Z)[-n]$: only $n=0,1$
survive.  The map (S09.51) then induces the identity
on $\mathbb Z\oplus\mathbb Z[-1]$; preservation of
the stated finite colimits and totalizations proves
the perfect-complex case.  These statements and
the complete proof were opened in the primary v2
source, Remarks~2.9 and~2.52, Proposition~3.6,
PDF pp.~13, 22, and~29--30
\cite{KubrakShuklinZakharov2026_check}.

There is a second explicit map in that paper which
is relevant to attempts to shift degrees before
using the deformation.  Rationalization
$M\to M\otimes\mathbb Q$ and the filtered monad map
$\mathbf L\operatorname{Sym}\to\mathbf L\operatorname{Bin}$
give
\[
 \alpha:\mathbf L\operatorname{Bin}(M)
       \longrightarrow
          \mathbf L\operatorname{Sym}_{\mathbb Z}
                              (M\otimes\mathbb Q).
 \tag{S09.52}
\]
Proposition~3.31 proves that this map is an
equivalence for $M\in D(\mathbb Z)^{\leq-1}$.
For $M=\mathbb Z[1]$, in particular, its positive
augmentation summand is $\mathbb Q[1]$.
Remark~3.34 explains the nonsplitting of the
filtration in this range; the associated graded
pieces are still the integral derived divided
powers in (S09.50).  These statements and the
proof of Proposition~3.31 were opened on PDF
pp.~39--40.  Section~1.1 explicitly reserves its
announced stronger deformation comparison on
$\mathbb Z[-3]$ for a sequel.  No such announced
comparison is used here.

The maps (S09.45), (S09.51), and (S09.52) provide
precise points of comparison for this route.
The obstruction (S09.49) proves why the direct
Rees lift of the prescribed integral coefficient
map is unavailable, while retaining the positive
derived equivalence (S09.51) in its full stated
range.  No separating invariant for all strict
admissible morphisms, and no common source for
arbitrary targets, is inferred from this audit.
